% !PRINTED 169
\weberpage{169}

%<<CONT
hervor, wenn man $\alpha=\beta=\frac12$, $\gamma=1$ setzt. Das eine ihrer
partikularen Integrale ist also
\weq{10}{F(\tfrac12,\tfrac12,1,x) = 1 + \sum_{\nu=1}^{\infty} \Bigl(\frac{1.3.5\ldots 2\nu-1}{2.4\ldots 2\nu}\Bigr)^2 x^\nu}
\[
{}= \sum_{\nu=0}^{\infty} \frac{\Pi(2\nu)^2}{\Pi(\nu)^4}\Bigl(\frac{x}{16}\Bigr)^\nu,
\]
und diese Reihe ist konvergent, so lange der absolute Wert von
$x$ kleiner als 1 ist. Als das zweite partikulare Integral kann man
wählen
\weq{11}{F(\tfrac12,\tfrac12,1,1-x),}
das aber einen anderen Konvergenzbereich hat. Denken wir uns
$x$ als komplexe Variable in einer Ebene dargestellt, so konvergiert
(10) in einem mit dem Radius 1 um den Nullpunkt beschriebenen
Kreise, (11) in einem gleichen Kreise um den Punkt 1 beschrieben,
so daß der gemeinsame Konvergenzbereich aus einem Zweieck
besteht, das den zwischen 0 und 1 gelegenen Teil der reellen
Achse enthält. Man kann aber auch, was für uns wichtig ist,
das zweite partikulare Integral in einer Form aufstellen, die in
demselben Kreise wie die Reihe (10) konvergiert. Dabei ist zu
beachten, daß das zweite partikulare Integral für $x=0$ unendlich
wird.

Dieses zweite partikulare Integral erhält man durch den
folgenden Grenzübergang: Bezeichnen wir die hypergeometrische
Reihe nach \gesperrt{Gauss} mit
\[
F(\alpha,\beta,\gamma,x) = 1 + \frac{\alpha.\beta}{1.\gamma}x + \frac{\alpha(\alpha+1)\beta(\beta+1)}{1.2.\gamma(\gamma+1)}x^2 + \cdots
\]
\[
{}= \frac{\Pi(\gamma-1)}{\Pi(\alpha-1)\Pi(\beta-1)} \sum_{\nu=0}^{\infty} \frac{\Pi(\alpha+\nu-1)\Pi(\beta+\nu-1)}{\Pi(\gamma+\nu-1)\Pi(\nu)} x^\nu,
\]
so erhält man im allgemeinen die beiden partikularen Integrale
von (9):
\[
\begin{aligned}
y_1 &= F_1(x) = F(\alpha,\beta,\gamma,x), \\
y_2 &= F_2(x) = x^{1-\gamma}(1-x)^{\gamma-\alpha-\beta} F(1-\alpha,1-\beta,2-\gamma,x).
\end{aligned}
\]
Da diese aber für $\alpha=\beta=\frac12$, $\gamma=1$ miteinander identisch
werden, so setze man zunächst $\alpha=\beta=\frac12$, $\gamma=1+\varepsilon$, also
\[
\begin{aligned}
y_1 &= F(\tfrac12,\tfrac12,1+\varepsilon,x), \\
y_2 &= \Bigl(\frac{1-x}{x}\Bigr)^\varepsilon F(\tfrac12,\tfrac12,1-\varepsilon,x).
\end{aligned}
\]
